\begin{initiativkarte}{SDS HAW – Sozialistisch. Demokratisch. Antifaschistisch.}

  % Bild (optional)
  \IfFileExists{bilder/sds_haw.jpg}{%
    \begin{center}
      \includegraphics[width=\linewidth,height=5cm,keepaspectratio]{bilder/sds_haw.jpg}
    \end{center}
    \vspace{0.4cm}
  }{}

  % Kurzbeschreibung
  \textit{\textcolor{hawgray}{Der SDS HAW ist eine linke Hochschulgruppe, die sich für soziale Gerechtigkeit, antifaschistische Hochschulpolitik und studentische Mitbestimmung einsetzt.}}

  \vspace{0.5cm}
  \hawrule

  % Beschreibung
  \textbf{\textcolor{hawblue}{Was machen wir?}}\\[0.3cm]

  Der SDS HAW ist eine sozialistisch-demokratische Hochschulgruppe an der HAW Hamburg. Wir setzen uns für eine solidarische, demokratische und kritische Hochschule ein.

  Unser Ziel ist eine Universität, die nicht nur auf den Arbeitsmarkt ausgerichtet ist, sondern gesellschaftliche Verantwortung übernimmt und allen Menschen gleiche Bildungschancen ermöglicht.

  Unsere Arbeit umfasst unter anderem:

  \begin{itemize}
      \item politische Hochschularbeit und Studierendenvertretung
      \item Organisation von Diskussionen, Veranstaltungen und Kampagnen
      \item Engagement gegen Sparzwänge im Bildungsbereich
      \item Einsatz für Antifaschismus und soziale Gerechtigkeit
      \item Beteiligung an akademischen Gremien und Wahlen
  \end{itemize}

  Der SDS HAW tritt regelmäßig bei Wahlen zum Studierendenparlament sowie in Fakultäts- und Departmenträten an.

  Ein zentraler Teil unserer Arbeit ist die politische Bildung und die gemeinsame Analyse gesellschaftlicher Strukturen im Hochschulkontext.

  Treffen finden regelmäßig montags um 18 Uhr in der Alexanderstraße 1 (Raum 0.14 EG) statt.

  \vspace{0.5cm}

  % Tags
  \tag{Politik}
  \tag{Studierendenvertretung}
  \tag{Antifaschismus}
  \tag{Gesellschaft}
  \tag{Diskussion}

  \vspace{0.6cm}
  \hawrule

  % Infozeilen
  \infozeile{\faUsers}{Zielgruppe: Politisch interessierte Studierende aller Fachrichtungen}
  \infozeile{\faGraduationCap}{Semester: Alle Semester willkommen}
  \infozeile{\faMapMarker*}{Ort: Alexanderstraße 1, Raum 0.14 (EG)}
  \infozeile{\faCalendarAlt}{Treffen: Montags, 18:00 Uhr}
  \infozeile{\faGlobe}{Website: \href{https://www.haw-hamburg.de}{SDS HAW Informationen}}

\end{initiativkarte}